% !TeX encoding = UTF-8
% !TeX spellcheck = en_US

% ====== 摘要 ======
\cabstract{
\thispagestyle{plain}
某某问题是……

本文\footnotecircle{本研究得到某某基金（编号：XXX）资助。}采用了……

研究表明……

使用说明：

1. 主文件为COEmain.tex，里面可详细看到对应封面，致谢，出版情况与个人简历对应的文件。

2. 默认指定使用图片在figures文件夹内，如现在使用的北大logo（logo.eps）、版权声明（copyright2.pdf）以及北京大学学位论文原创性声明和使用授权说明（Authority2.pdf）。

3. 参考文献使用基于BIB文件管理，建议对应的BIB代码使用学术谷歌统一标定。

4. 编译软件大家可以选择使用CTEX或TEXLIVE,需要保证完全安装。编译过程是一次xelatex命令后，再文献编译bibtex一次，再xelatex两次，即可查看pdf结果。

5. 当论文中包含大量eps格式的图片时，xelatex编译时间较长；将所有eps图片的格式转为pdf可以显著缩短编译时间。

6. 打印生成的pdf时，为了保证与学院Word模板相同的效果，请选择“实际大小”或者“自定义比例100\%”。
\smallskip

{\color{red}{英文培养项目学生使用英文写学位论文，中文摘要应不少于6000字。}}

{\color{red}{Students of the English Training Program are required to write their dissertations in English, and their Chinese abstracts should be no less than 6,000 Chinese characters.}}
\smallskip
%\ \quad
% 当关键词与脚注（某某基金资助）在同一页时，需要人工增减“\quad \\”来保证“关键词放摘要页最下方”。另一种方法是采用\vfill，但是在package.tex中使用脚注控宏包footmisc时就不能再采用选项bottom，代价是每页的脚注底部可能不对齐。

% 当中文摘要较长，使得关键词在摘要的第二页时，采用\vfill即可使关键词位于页面底部。（与英文摘要中关键词的处理方法相同）
 
\vspace{7.1cm} 关键词：关键词1，关键词2，关键词3……{\color{red}{关键词需要置底}}
}


% ====== Abstract ======
\eabstract{
\thispagestyle{plain}

In environmental economics, environmental resources including environmental quality are categorized as amenity resources. Due to its importance to human welfare, the amenity resources theoretical study and valuation is an ongoing issue at the academic frontier in the environmental economics area.


\vfill

KEYWORDS: Key word 1, Key word 2, Key word 3, ……
}
